\textbf{COLUNI 2020 -- Questão 09}\par\medskip Leia o enunciado que segue.

\par

“A justiça \underline{estava} lá, e o meu máximo não \underline{fora} o suficiente.” CARRASCOZA, J. A. \textit{Aos 7 e aos 40}. Rio de Janeiro: Alfaguara, 2016. p. 58.

\par

Considere as afirmativas que apresentam a mesma sequência de tempos e de modos dos verbos, no exemplo acima, atribuindo V para a(s) verdadeira(s) e F para a(s) falsa(s).

\par

( ) “[...] é preciso mesmo olhar pra trás se queremos ir em frente.” (CARRASCOZA, J. A. \textit{Aos 7 e aos 40}. Rio de Janeiro: Alfaguara, 2016. p. 59.) ( ) “[...] quando eu já nem lembrava mais por que tínhamos viajado até ali [...].” (CARRASCOZA, J. A. \textit{Aos 7 e aos 40}. Rio de Janeiro: Alfaguara, 2016. p. 95.) ( ) “[...] Eu senti que ele se envergonharia se eu percebesse.” (CARRASCOZA, J. A. \textit{Aos 7 e aos 40}. Rio de Janeiro: Alfaguara, 2016. p. 39.) ( ) “Sim, o garoto precisava de cuidados, o médico disse, haviam agido bem, [...]” (CARRASCOZA, J. A. \textit{Aos 7 e aos 40}. Rio de Janeiro: Alfaguara, 2016. p. 32.)

\par

Assinale a sequência CORRETA:\begin{enumerate}
\item V, V, V, V.
\item V, F, V, F.
\item F, V, F, F.
\item F, V, F, V.\end{enumerate}
